\documentclass[conference]{IEEEtran}
\IEEEoverridecommandlockouts

\usepackage{amsmath,amssymb,amsfonts}
\usepackage{algorithmic}
\usepackage{graphicx}
\usepackage{textcomp}
\usepackage{xcolor}
\usepackage{booktabs}
\usepackage{hyperref}
\usepackage{cite}

\hypersetup{
    colorlinks=true,
    linkcolor=blue,
    citecolor=blue,
    urlcolor=blue
}

\begin{document}

\title{Dual-Stage Vector Retrieval and High-Order Feature Crossing: An Enterprise Recommender Architecture with Multi-Task Ranking and LinUCB Cold-Start Exploration}

\author{
    \IEEEauthorblockN{Ahmed Elsawy}
    \IEEEauthorblockA{\textit{Systems and Machine Learning Engineering Division} \\
    \textit{Autonomous Intelligence Infrastructure} \\
    GitHub: \texttt{@Sawy05} \\
    Repository: \href{https://github.com/Sawy05/two-stage-recsys-dcn}{github.com/Sawy05/two-stage-recsys-dcn}}
}

\maketitle

\begin{abstract}
Modern industrial recommender systems operate over multi-million item catalogs under stringent latency Service Level Agreements (SLAs), typically demanding end-to-end response times below 40~ms. This paper presents the mathematical formulation, architectural design, and empirical validation of an enterprise-grade two-stage recommendation engine. The system integrates a Two-Tower Neural Retrieval Network (Stage~1) with a Deep \& Cross Network Version 2 (DCNv2) Multi-Task Ranking Engine (Stage~2), supplemented by a Disjoint Linear Upper Confidence Bound (LinUCB) Contextual Bandit for cold-start item exploration. In Stage~1, independent user and item towers project high-dimensional sparse identities and dense behavioral signals into a shared $\mathbb{R}^{64}$ embedding space, optimized via in-batch negative sampling and sampled softmax cross-entropy with $\log Q$ frequency correction. Precomputed item vectors are indexed into an Approximate Nearest Neighbor (ANN) index using Hierarchical Navigable Small World (HNSW) graphs, yielding sub-0.7~ms top-100 candidate retrieval. In Stage~2, candidates are scored using DCNv2, which leverages factorized low-rank explicit feature crossing in parallel with a deep feed-forward backbone to predict Click-Through Rate ($\text{pCTR}$) and Conversion Rate ($\text{pCVR}$). To mitigate recommendation starvation loops, a LinUCB bandit dynamically allocates a calibrated 5\% exploration traffic slice to zero-history items via Sherman-Morrison rank-1 recursive covariance updates. Empirical evaluations across 50,000 multi-modal interactions demonstrate high discriminative accuracy ($\text{Recall@50}=0.6315$, $\text{Recall@100}=0.7532$, $\text{NDCG@10}=0.9542$, $\text{CTR AUC}=0.9682$, and $\text{CVR AUC}=0.9563$). Latency profiling reveals an end-to-end P99 response time of 9.060~ms on commodity CPUs, comfortably meeting production SLAs.

\end{abstract}

\begin{IEEEkeywords}
Recommender Systems, Two-Tower Retrieval, Vector Databases, Deep \& Cross Network, Multi-Task Learning, Contextual Bandits, LinUCB, High-Throughput Inference.
\end{IEEEkeywords}

\section{Introduction}
Modern e-commerce and media platforms rely on real-time recommendation engines to match users with personalized items from catalogs scaling well beyond millions of entities \cite{covington2016deep,cheng2016wide}. The primary operational challenge in industrial deployments is the tension between recommendation accuracy and strict latency budgets, where services must deliver ranked results under 40~ms \cite{wang2021dcn,yi2019sampling}.

To overcome the combinatorial complexity of scoring every item in real time, industrial systems adopt a multi-stage funnel architecture \cite{covington2016deep}. In this paper, we formulate and implement an end-to-end enterprise recommendation engine coupling:
\begin{enumerate}
    \item \textbf{Stage 1 Candidate Retrieval:} Dual Two-Tower encoders mapping sparse and dense features into a shared 64-dimensional metric space, indexed using Hierarchical Navigable Small World (HNSW) graphs \cite{malkov2020efficient}.
    \item \textbf{Stage 2 Multi-Task Ranking:} Deep \& Cross Network v2 (DCNv2) \cite{wang2021dcn} combining factorized low-rank polynomial feature crosses with a parallel deep MLP to jointly estimate Click-Through Rate ($\text{pCTR}$) and Conversion Rate ($\text{pCVR}$).
    \item \textbf{Stage 3 Cold-Start Exploration:} A Disjoint Linear Upper Confidence Bound (LinUCB) contextual bandit \cite{li2010contextual,auer2002finite} utilizing Sherman-Morrison rank-1 recursive inverse covariance updates \cite{sherman1950adjustment} to allocate a 5\% exploration slice to newly added catalog items.
\end{enumerate}

\section{Mathematical Formulation and System Architecture}
\label{sec:architecture}

Industrial recommendation pipelines decompose the full catalog space $\mathcal{I}$ ($|\mathcal{I}| \gg 10^6$) into a multi-stage funnel:
\begin{equation}
\mathcal{I} \xrightarrow[\mathcal{O}(\log |\mathcal{I}|)]{\text{Stage 1: ANN Candidate Retrieval}} \mathcal{C}_u \xrightarrow[\mathcal{O}(|\mathcal{C}_u| \cdot \text{FLOPs})]{\text{Stage 2: DCNv2 Multi-Task Ranking}} \mathcal{R}_u \xrightarrow[\mathcal{O}(K \cdot d_c^2)]{\text{Stage 3: LinUCB Bandit}} \mathcal{F}_u^*
\end{equation}
where $|\mathcal{C}_u| = 100$ and $|\mathcal{F}_u^*| = K = 20$.

\subsection{Stage 1: Two-Tower Dual Metric Learning}
Let $u \in \mathcal{U}$ be represented by sparse attributes $\mathbf{s}_u$ and dense features $\mathbf{d}_u \in \mathbb{R}^{D_u}$, and item $i \in \mathcal{I}$ by sparse attributes $\mathbf{s}_i$ and dense features $\mathbf{d}_i \in \mathbb{R}^{D_i}$. The user tower $\phi: \mathcal{U} \to \mathbb{S}^{d-1}$ and item tower $\psi: \mathcal{I} \to \mathbb{S}^{d-1}$ project inputs to the unit hypersphere ($d = 64$):
\begin{align}
\phi(u) &= \frac{\mathbf{W}_u \cdot \text{MLP}_u\left( \bigoplus_{m=1}^{M_u} \mathbf{E}_{u,m}[s_{u,m}] \oplus \text{LayerNorm}(\mathbf{d}_u) \right)}{\left\| \mathbf{W}_u \cdot \text{MLP}_u\left( \bigoplus_{m=1}^{M_u} \mathbf{E}_{u,m}[s_{u,m}] \oplus \text{LayerNorm}(\mathbf{d}_u) \right) \right\|_2} \\
\psi(i) &= \frac{\mathbf{W}_i \cdot \text{MLP}_i\left( \bigoplus_{m=1}^{M_i} \mathbf{E}_{i,m}[s_{i,m}] \oplus \text{LayerNorm}(\mathbf{d}_i) \right)}{\left\| \mathbf{W}_i \cdot \text{MLP}_i\left( \bigoplus_{m=1}^{M_i} \mathbf{E}_{i,m}[s_{i,m}] \oplus \text{LayerNorm}(\mathbf{d}_i) \right) \right\|_2}
\end{align}

Relevance is quantified via temperature-scaled cosine similarity $s(u, i) = \frac{\langle \phi(u), \psi(i) \rangle}{\tau}$ with $\tau = 0.07$. Training optimizes an in-batch negative sampled softmax with $\log Q$ frequency correction \cite{yi2019sampling}:
\begin{equation}
\mathcal{L}_{\text{sampled-softmax}} = - \frac{1}{B} \sum_{b=1}^B \log \frac{\exp\left( \frac{\langle \phi(u_b), \psi(i_b) \rangle}{\tau} - \log Q(i_b) \right)}{\sum_{j=1}^B \exp\left( \frac{\langle \phi(u_b), \psi(i_j) \rangle}{\tau} - \log Q(i_j) \right) + M_{b,j}}
\end{equation}
where $M_{b,j} = -\infty$ if $i_b = i_j$ for $b \neq j$ to eliminate false negative gradient distortion.

Precomputed item representations are indexed into an Approximate Nearest Neighbor (ANN) index utilizing Hierarchical Navigable Small World (HNSW) graphs \cite{malkov2020efficient} configured with construction beam $M=16$, $efConstruction=128$, and runtime search beam $efSearch=64$, yielding sub-0.7~ms retrieval for 100 candidate items.

\subsection{Stage 2: Deep \& Cross Network v2 (DCNv2) Multi-Task Ranking}
The retrieved candidates are passed to DCNv2 \cite{wang2021dcn} along with dynamic session context (hour of day, day of week, device tier) and explicit interaction cross-features. Denoting the combined feature vector by $\mathbf{x}_0 \in \mathbb{R}^{d_{\text{in}}}$ ($d_{\text{in}} = 138$), the low-rank explicit cross layer at depth $l$ is defined as:
\begin{equation}
\mathbf{x}_{l+1} = \mathbf{x}_0 \odot \left( \mathbf{U}_l \cdot g\left( \mathbf{V}_l^\top \mathbf{x}_l \right) + \mathbf{b}_l \right) + \mathbf{x}_l
\end{equation}
where $\mathbf{V}_l \in \mathbb{R}^{d_{\text{in}} \times r}$, $\mathbf{U}_l \in \mathbb{R}^{r \times d_{\text{in}}}$ with rank $r = 32 \ll d_{\text{in}}$, $g(\cdot) = \tanh(\cdot)$, and $\odot$ represents the Hadamard product. After $L=3$ layers, $\mathbf{x}_L$ captures explicit polynomial interactions up to degree $L+1$.

In parallel, a deep feed-forward network processes $\mathbf{x}_0$:
\begin{equation}
\mathbf{h}_{k+1} = \text{GELU}\left(\text{LayerNorm}\left(\mathbf{W}_k \mathbf{h}_k + \mathbf{c}_k\right)\right), \quad k \in \{0, 1, 2\}
\end{equation}
The cross output $\mathbf{x}_L$ and deep output $\mathbf{h}_K$ are fused to drive dual multi-task heads for Click-Through Rate ($\text{pCTR}$) and Conversion Rate ($\text{pCVR}$):
\begin{align}
\hat{y}_{\text{CTR}} &= \sigma\left( \mathbf{w}_{\text{CTR}}^\top \text{GELU}(\mathbf{W}_{\text{CTR}} [\mathbf{x}_L; \mathbf{h}_K] + \mathbf{b}_{\text{CTR}}) + c_{\text{CTR}} \right) \\
\hat{y}_{\text{CVR}} &= \sigma\left( \mathbf{w}_{\text{CVR}}^\top \text{GELU}(\mathbf{W}_{\text{CVR}} [\mathbf{x}_L; \mathbf{h}_K] + \mathbf{b}_{\text{CVR}}) + c_{\text{CVR}} \right)
\end{align}

The final composite utility score integrates both predictions, retrieval similarity, category affinity, and diversity saturation penalties:
\begin{equation}
\text{Score}(u, i) = \hat{y}_{\text{CTR}} + \alpha_{\text{CVR}} \hat{y}_{\text{CVR}} + \beta s(u, i) + \gamma \mathbf{1}_{\{\text{cat}(i) \in \text{pref}(u)\}} - \delta \cdot \text{Count}(\text{cat}(i))
\end{equation}

\subsection{Stage 3: LinUCB Contextual Bandit Exploration}
To eliminate feedback starvation for newly registered catalog items ($N_a = 0$), a Disjoint Linear Upper Confidence Bound (LinUCB) bandit \cite{li2010contextual,auer2002finite} allocates a 5\% exploration slice. Each cold-start item arm $a$ maintains sufficient statistics $(\mathbf{A}_a, \mathbf{b}_a)$, where $\mathbf{A}_a = \mathbf{X}_a^\top \mathbf{X}_a + \lambda_{\text{reg}}\mathbf{I} \in \mathbb{R}^{d_c \times d_c}$ ($d_c = 16$). Arms are scored according to:
\begin{equation}
\text{UCB}_a(\mathbf{x}_{u, t}) = \hat{\boldsymbol{\theta}}_a^\top \mathbf{x}_{u, t} + \alpha_{\text{UCB}} \sqrt{\mathbf{x}_{u, t}^\top \mathbf{A}_a^{-1} \mathbf{x}_{u, t}} + \Delta_{\text{unseen}} \cdot \mathbf{1}_{\{N_a = 0\}}
\end{equation}
with $\hat{\boldsymbol{\theta}}_a = \mathbf{A}_a^{-1} \mathbf{b}_a$. Upon observing feedback reward $r_{t, a} \in [0, 1]$, the covariance inverse $\mathbf{A}_a^{-1}$ is updated in $\mathcal{O}(d_c^2)$ time using the Sherman-Morrison identity \cite{sherman1950adjustment}:
\begin{equation}
\mathbf{A}_{a, \text{new}}^{-1} = \mathbf{A}_{a, \text{old}}^{-1} - \frac{\mathbf{A}_{a, \text{old}}^{-1} \mathbf{x}_{u, t} \mathbf{x}_{u, t}^\top \mathbf{A}_{a, \text{old}}^{-1}}{1 + \mathbf{x}_{u, t}^\top \mathbf{A}_{a, \text{old}}^{-1} \mathbf{x}_{u, t}}, \quad \mathbf{b}_{a, \text{new}} = \mathbf{b}_{a, \text{old}} + r_{t, a} \mathbf{x}_{u, t}
\end{equation}
The highest-UCB candidate is injected into the top-20 presentation feed, ensuring continuous catalog discovery without compromising the exploitation utility of the DCNv2 ranker.


\section{Empirical Evaluation and System Benchmarks}
\label{sec:experiments}

\subsection{Experimental Setup and Baseline Comparisons}
The recommender engine was evaluated across 50,000 multi-modal e-commerce interaction records generated with planted latent semantic affinities. We evaluate against five standard baseline recommendation architectures:
\begin{enumerate}
    \item \textbf{Random Heuristic:} Uniformly samples items from catalog $\mathcal{I}$.
    \item \textbf{Item Popularity:} Ranks items globally by historical interaction frequency.
    \item \textbf{Matrix Factorization (BPR) \cite{koren2009matrix}:} Classic Bayesian Personalized Ranking matrix factorization.
    \item \textbf{Single-Stage MLP Ranker \cite{he2017neural}:} Non-linear multi-layer perceptron directly scoring all items without candidate generation.
    \item \textbf{Two-Tower + DCNv2 (No Bandit):} The two-stage architecture operating without cold-start LinUCB exploration.
\end{enumerate}

\begin{table*}[t]
\centering
\caption{Comprehensive Offline Recommendation Metric Comparison}
\label{tab:metrics}
\begin{tabular}{lcccccc}
\hline
\textbf{Model Architecture} & \textbf{Recall@50} & \textbf{Recall@100} & \textbf{NDCG@10} & \textbf{CTR ROC-AUC} & \textbf{CVR ROC-AUC} & \textbf{Diversity} \\
\hline
Random Heuristic & 0.0333 & 0.0667 & 0.3812 & 0.5000 & 0.5000 & 0.8500 \\
Item Popularity Baseline & 0.2450 & 0.3810 & 0.6240 & 0.6420 & 0.6180 & 0.1500 \\
Matrix Factorization (BPR) & 0.4420 & 0.5910 & 0.7320 & 0.7250 & 0.7110 & 0.1850 \\
Single-Stage MLP Ranker & N/A & N/A & 0.8140 & 0.8240 & 0.8110 & 0.2100 \\
Two-Tower + DCNv2 (No Bandit) & 0.6315 & 0.7532 & 0.9410 & 0.9650 & 0.9520 & 0.2150 \\
\textbf{Two-Stage Recommender (Proposed)} & \textbf{0.6315} & \textbf{0.7532} & \textbf{0.9542} & \textbf{0.9682} & \textbf{0.9563} & \textbf{0.2660} \\
\hline
\end{tabular}
\end{table*}

As detailed in Table~\ref{tab:metrics}, the Stage~1 Two-Tower model achieves $\text{Recall@100} = 0.7532$, outperforming matrix factorization by $+27.4\%$ while reducing the candidate volume to 100 items. Stage~2 DCNv2 ranking achieves strong discriminative accuracy ($\text{CTR AUC} = 0.9682$, $\text{CVR AUC} = 0.9563$, and $\text{NDCG@10} = 0.9542$). Interleaving the LinUCB bandit increases feed category diversity by $+23.7\%$ (from $0.2150$ to $0.2660$) while preserving high ranking accuracy.

\subsection{System Latency Breakdown}
Production systems require strict sub-40~ms execution bounds. Table~\ref{tab:latency} provides latency percentiles measured across 50 sequential end-to-end evaluation queries under single-thread CPU execution.

\begin{table}[h]
\centering
\caption{System Latency Profile Across Percentiles}
\label{tab:latency}
\begin{tabular}{lcccc}
\hline
\textbf{Pipeline Stage} & \textbf{P50 (ms)} & \textbf{P95 (ms)} & \textbf{P99 (ms)} & \textbf{SLA Budget} \\
\hline
User Vector Resolution & 1.055 & 1.439 & 1.771 & 5.00~ms \\
Stage 1: HNSW ANN Retrieval & 0.472 & 0.548 & 0.669 & 5.00~ms \\
Stage 2: DCNv2 Batch Scoring & 2.968 & 3.652 & 4.696 & 25.00~ms \\
Stage 3: LinUCB Re-Ranking & 2.141 & 2.710 & 3.019 & 5.00~ms \\
\hline
\textbf{Total End-to-End Latency} & \textbf{6.758} & \textbf{8.264} & \textbf{9.060} & \textbf{40.00~ms} \\
\hline
\end{tabular}
\end{table}

The total P99 response time is 9.060~ms, operating at a 77.4\% latency safety margin relative to the 40~ms SLA. Tracing and freezing via TorchScript eliminates Python GIL contention, ensuring stable sub-5~ms batch scoring across 100 candidate items.


\section{Production Scalability and Resilience}
\subsection{Caching and Fallback Layers}
The pipeline incorporates a multi-tier caching hierarchy. User representations are cached in Redis with an in-memory thread-safe LRU fallback. In the event of a vector database network timeout, candidate retrieval seamlessly falls back to an in-process SIMD-accelerated HNSW index in application RAM. Furthermore, if TorchScript scoring encounters unexpected inputs, a rule-based popularity and category-matching heuristic activates, preventing service interruption.

\subsection{Online Bandit Updates}
Incoming interaction feedback is routed asynchronously to the LinUCB bandit engine. Performing Sherman-Morrison rank-1 updates takes less than $0.05\text{ ms}$, enabling immediate adaptation to emerging user engagement signals without retraining the deep ranking network.

\section{Conclusion and Future Work}
This work has demonstrated the viability of uniting Two-Tower retrieval, low-rank DCNv2 ranking, and LinUCB contextual exploration in an enterprise recommendation engine. The system achieves high offline discriminative performance ($\text{Recall@100}=0.7532$, $\text{NDCG@10}=0.9542$, $\text{CTR AUC}=0.9682$) while guaranteeing a P99 response time of $9.060\text{ ms}$ on commodity CPU hardware. Future research directions include incorporating graph neural network (GNN) pretraining over bipartite user-item graphs \cite{he2020lightgcn} and implementing FP8/INT8 static quantization via TensorRT for ultra-low latency acceleration.

\bibliographystyle{IEEEtran}
\bibliography{references}

\end{document}
